\documentclass[a4paper,12pt,titlepage,finall]{article}

\usepackage[T1,T2A]{fontenc}     % форматы шрифтов
\usepackage[utf8x]{inputenc}     % кодировка символов, используемая в данном файле
\usepackage[russian]{babel}      % пакет русификации
\usepackage{tikz}                % для создания иллюстраций
\usepackage{pgfplots}            % для вывода графиков функций
\usepackage{geometry}            % для настройки размера полей
\usepackage{indentfirst}         % для отступа в первом абзаце секции
\usepackage{amsmath}

% выбираем размер листа А4, все поля ставим по 3см
\geometry{a4paper,left=30mm,top=30mm,bottom=30mm,right=30mm}

\setcounter{secnumdepth}{0}      % отключаем нумерацию секций

\pgfplotsset{compat=1.18}
\usetikzlibrary{arrows.meta,positioning}
\usepgfplotslibrary{fillbetween} % для изображения областей на графиках

\begin{document}
% Титульный лист
\begin{titlepage}
    \begin{center}
    {\small \sc Московский государственный университет \\имени М.~В.~Ломоносова\\
    Факультет вычислительной математики и кибернетики\\}
    \vfill
    {\Large \sc Отчет по заданию №6}\\
    ~\\
    {\large \bf <<Сборка многомодульных программ. \\
    Вычисление корней уравнений и определенных интегралов.>>}\\
    ~\\
    {\large \bf Вариант 8}
    \end{center}
    \begin{flushright}
    \vfill {Выполнил:\\
    студент 104 группы\\
    Д.~М.~Алексеевич\\
    ~\\
    Преподаватель:\\
    Д.~А.~Гуляев
    }
    \end{flushright}
    \begin{center}
    \vfill
    {\small Москва\\2026}
    \end{center}
\end{titlepage}

% Автоматически генерируем оглавление на отдельной странице
\tableofcontents
\newpage

\section{Постановка задачи}

Требуется реализовать многомодульную программу для вычисления площади
плоской фигуры, ограниченной тремя кривыми
\[
    y=f_1(x),\qquad y=f_2(x),\qquad y=f_3(x).
\]
Для варианта 8 заданы функции
\[
    f_1(x)=e^x+2,\qquad
    f_2(x)=-2x+8,\qquad
    f_3(x)=-\frac{5}{x}.
\]

Абсциссы вершин фигуры находятся как корни уравнений $f_i(x)-f_j(x)=0$.
В программе реализованы два метода приближенного решения уравнений:
метод Ньютона и метод хорд. Выбор метода выполняется во время сборки
с помощью макроса препроцессора. Определенные интегралы вычисляются
составной формулой Симпсона.

Функции $f_1$, $f_2$, $f_3$ и их производные реализованы на ассемблере
NASM32 с использованием арифметики x87 и соглашения вызова cdecl.
Остальная часть программы написана на языке C. Тип вычислений:
\texttt{long double}.

\newpage

\section{Математическое обоснование}

Графики заданных функций показаны на рис.~\ref{plot1}. Функция $f_3$
имеет разрыв в точке $x=0$, поэтому при выборе отрезков локализации
корней необходимо рассматривать левую и правую полуплоскости отдельно.
Анализ знаков разностей функций показывает, что замкнутая область
образуется точками пересечения $f_1=f_3$, $f_2=f_3$ и $f_1=f_2$.

\begin{figure}[h]
\centering
\begin{tikzpicture}
\begin{axis}[
             axis lines=middle,
             restrict x to domain=-3.5:2.5,
             restrict y to domain=-1:12,
             enlargelimits,
             legend cell align=left,
             scale=1.2]

\addplot[green,domain=-3.5:2.5,samples=256,thick] {exp(x)+2};
\addlegendentry{$y=e^x+2$}

\addlegendimage{empty legend}\addlegendentry{}

\addplot[blue,domain=-3.5:2.5,samples=256,thick] {-2*x+8};
\addlegendentry{$y=-2x+8$}

\addplot[red,domain=-3.5:-0.08,samples=256,thick] {-5/x};
\addlegendentry{$y=-\frac{5}{x}$}

\end{axis}
\end{tikzpicture}
\caption{Графики функций варианта 8}
\label{plot1}
\end{figure}

Для поиска корней выбраны следующие отрезки:
\[
    x_{13}\in[-3,-2],\qquad
    x_{23}\in[-1,-0.1],\qquad
    x_{12}\in[1,2].
\]
На концах каждого из этих отрезков соответствующая разность функций
меняет знак, а сами функции непрерывны на выбранных отрезках. Поэтому
на каждом отрезке существует корень. Условия применимости используемых
численных методов и оценки погрешности вычисления корней и интегралов
следуют из стандартных результатов математического анализа~\cite{math}.

Для достижения точности вычисления площади порядка $10^{-4}$ выбраны
более строгие внутренние точности
\[
    \varepsilon_1=10^{-8},\qquad \varepsilon_2=10^{-8}.
\]
При интегрировании используется правило Рунге для формулы Симпсона:
\[
    \frac{|I_{2n}-I_n|}{15}<\varepsilon_2.
\]

\newpage

\section{Результаты экспериментов}

Координаты найденных точек пересечения приведены в таблице~\ref{table1}.

\begin{table}[h]
\centering
\begin{tabular}{|c|c|c|}
\hline
Кривые & $x$ & $y$ \\
\hline
1 и 3 & -2.390536703909 & 2.091580519063 \\
2 и 3 & -0.549509756796 & 9.099019513596 \\
1 и 2 &  1.251757931391 & 5.496484137217 \\
\hline
\end{tabular}
\caption{Координаты точек пересечения}
\label{table1}
\end{table}

Площадь фигуры представляется как сумма двух интегралов:
\[
S =
\int_{x_{13}}^{x_{23}}\bigl(f_3(x)-f_1(x)\bigr)\,dx+
\int_{x_{23}}^{x_{12}}\bigl(f_2(x)-f_1(x)\bigr)\,dx.
\]
В результате вычислений получено
\[
    S\approx 9.806944933573.
\]

\begin{figure}[h]
\centering
\begin{tikzpicture}
\begin{axis}[
             axis lines=middle,
             restrict x to domain=-3.5:2.5,
             restrict y to domain=-1:12,
             enlargelimits,
             legend cell align=left,
             scale=1.2,
             xticklabels={,,},
             yticklabels={,,}]

\addplot[green,domain=-3.5:2.5,samples=256,thick,name path=A] {exp(x)+2};
\addlegendentry{$y=e^x+2$}

\addlegendimage{empty legend}\addlegendentry{}

\addplot[blue,domain=-3.5:2.5,samples=256,thick,name path=B] {-2*x+8};
\addlegendentry{$y=-2x+8$}

\addplot[red,domain=-3.5:-0.08,samples=256,thick,name path=C] {-5/x};
\addlegendentry{$y=-\frac{5}{x}$}

\addplot[blue!20,samples=256] fill between[of=C and A,soft clip={domain=-2.390536703909:-0.549509756796}];
\addplot[blue!20,samples=256] fill between[of=B and A,soft clip={domain=-0.549509756796:1.251757931391}];
\addlegendentry{$S=9.8069$}

\addplot[dashed] coordinates { (-2.390536703909, 2.091580519063) (-2.390536703909, 0) };
\addplot[color=black] coordinates {(-2.390536703909, 0)} node [label={-135:{\small -2.3905}}]{};

\addplot[dashed] coordinates { (-0.549509756796, 9.099019513596) (-0.549509756796, 0) };
\addplot[color=black] coordinates {(-0.549509756796, 0)} node [label={-90:{\small -0.5495}}]{};

\addplot[dashed] coordinates { (1.251757931391, 5.496484137217) (1.251757931391, 0) };
\addplot[color=black] coordinates {(1.251757931391, 0)} node [label={-45:{\small 1.2518}}]{};

\end{axis}
\end{tikzpicture}
\caption{Плоская фигура, ограниченная графиками заданных уравнений}
\label{plot2}
\end{figure}

\newpage

\section{Структура программы и спецификация функций}

Программа состоит из следующих модулей:
\begin{itemize}
\item \texttt{src/main.c} --- основной модуль, содержащий функции
    \texttt{root}, \texttt{integral}, обработку ключей командной строки,
    тесты и вычисление площади;
\item \texttt{src/functions.asm} --- ассемблерный модуль, содержащий
    функции \texttt{f1}, \texttt{f2}, \texttt{f3}, \texttt{df1},
    \texttt{df2}, \texttt{df3};
\item \texttt{Makefile} --- сценарий сборки программы.
\end{itemize}

Основные функции программы:
\begin{itemize}
\item \texttt{root} --- вычисляет корень уравнения $f(x)=g(x)$ на
    заданном отрезке;
\item \texttt{integral} --- вычисляет определенный интеграл функции на
    заданном отрезке по формуле Симпсона;
\item \texttt{find\_roots} --- находит три вершины фигуры;
\item \texttt{find\_area} --- вычисляет площадь по найденным вершинам;
\item \texttt{test\_root}, \texttt{test\_integral} --- выполняют
    тестирование численных методов.
\end{itemize}

\begin{figure}[h]
\centering
\begin{tikzpicture}[node distance=17mm, every node/.style={draw, rounded corners, align=center}]
\node (main) {main.c\\CLI и управление};
\node (root) [below left of=main, xshift=-18mm] {root\\Ньютон/хорды};
\node (integral) [below right of=main, xshift=18mm] {integral\\Симпсон};
\node (asm) [below of=root, yshift=-5mm] {functions.asm\\$f_i$, $f_i'$};
\node (area) [below of=integral, yshift=-5mm] {вычисление\\площади};
\draw[->] (main) -- (root);
\draw[->] (main) -- (integral);
\draw[->] (root) -- (asm);
\draw[->] (integral) -- (asm);
\draw[->] (integral) -- (area);
\end{tikzpicture}
\caption{Структура программы}
\label{structure}
\end{figure}

\newpage

\section{Сборка программы (Make-файл)}

Сборка выполняется утилитой \texttt{make}. Цель \texttt{all} собирает
исполняемый файл \texttt{lab6}. Цель \texttt{chord} пересобирает
программу с макросом \texttt{USE\_CHORD}, включая метод хорд вместо
метода Ньютона. Цель \texttt{clean} удаляет объектные файлы и
исполняемый файл.

\begin{verbatim}
CC = gcc
ASM = nasm

TARGET = lab6
SRC_DIR = src
BUILD_DIR = build

CFLAGS = -m32 -Wall -Wextra -std=c99 -O2
LDFLAGS = -m32
LDLIBS = -lm

METHOD ?= newton

ifeq ($(METHOD),chord)
    CFLAGS += -DUSE_CHORD
endif
\end{verbatim}

Зависимости модулей:
\[
    \texttt{lab6} \leftarrow
    \texttt{main.o}+\texttt{functions.o},
\]
\[
    \texttt{main.o} \leftarrow \texttt{src/main.c},\qquad
    \texttt{functions.o} \leftarrow \texttt{src/functions.asm}.
\]

\newpage

\section{Отладка программы, тестирование функций}

Для проверки функции \texttt{root} используется уравнение
\[
    f_2(x)=f_3(x).
\]
Оно сводится к квадратному уравнению
\[
    -2x+8=-\frac{5}{x},\qquad 2x^2-8x-5=0.
\]
Корень, лежащий на выбранном отрезке, равен
\[
    x=\frac{4-\sqrt{26}}{2}\approx -0.549509756796.
\]
Этот результат используется в режиме \texttt{--test-root}.

Для проверки функции \texttt{integral} используется интеграл линейной
функции:
\[
    \int_0^1(-2x+8)\,dx = 7.
\]
Этот результат используется в режиме \texttt{--test-integral}.
Все тесты запускаются ключом \texttt{--test} или целью \texttt{make test}.

\newpage

\section{Программа на Си и на Ассемблере}

Исходные тексты программы находятся в каталоге \texttt{src}. Основная
логика численных методов реализована в файле \texttt{main.c}. Функции
варианта и их производные реализованы в файле \texttt{functions.asm}.
В ассемблерном модуле используются инструкции x87; результат каждой
функции возвращается в регистре \texttt{st0}, что соответствует
соглашению возврата вещественного значения в 32-разрядном режиме.

Программа поддерживает следующие ключи командной строки:
\begin{itemize}
\item \texttt{--help} --- справка;
\item \texttt{--roots} --- вывод точек пересечения;
\item \texttt{--iterations} --- вывод числа итераций;
\item \texttt{--test-root} --- тестирование функции \texttt{root};
\item \texttt{--test-integral} --- тестирование функции \texttt{integral};
\item \texttt{--test} --- запуск всех тестов.
\end{itemize}

\newpage

\section{Анализ допущенных ошибок}

Основные источники погрешности связаны с приближенным вычислением корней,
численным интегрированием и округлением вещественных чисел. Для уменьшения
влияния этих ошибок внутренние точности $\varepsilon_1$ и $\varepsilon_2$
выбраны существенно меньше требуемой точности площади. При вычислении
интегралов используется уточнение по правилу Рунге, а при поиске корней
используются отрезки, на которых соответствующие разности функций
непрерывны и меняют знак.

\newpage
\begin{raggedright}
\addcontentsline{toc}{section}{Список цитируемой литературы}
\begin{thebibliography}{99}
\bibitem{math} Ильин~В.\,А., Садовничий~В.\,А., Сендов~Бл.\,Х. Математический анализ. Т.\,1~---
    Москва: Наука, 1985.
\end{thebibliography}
\end{raggedright}

\end{document}
